\newsec{Tensione di Hall al variare della corrente}{Stima di $ R_H $ valutando l'andamento di $V_H(i)$.}

Per sfruttare l'\cref{eq:hall-tension} abbiamo misurato la tensione di Hall al variare della corrente $ i_p $ per $ 10 $ valori di campo magnetico $ B $ diversi.

\subsection{Contributo disallineamento piastre}

Per $ B=0 $ ci si aspetta che $ V_H \equiv 0 $ per qualsiasi valore di corrente.
In realtà, sperimentalmente, si osserva (\cref{fig:ohmic-contribution}) un andamento lineare tra la tensione di Hall e la corrente iniettata.

Questo contributo è dovuto al disallineamento tra le facce del campione di germanio: siamo interessati a caratterizzarlo per sottrarlo all'andamento di $ V_H(i_p) $ dovuto all'effetto Hall.

\subsection{Regressioni su $ V_{H}(i_p) $}

Per ciascuno dei $ 10 $ set di dati di raccolti abbiamo effettuato una regressione lineare:
\begin{equation}
V_H = q + m \cdot i_p
\end{equation}
aspettandoci, per confronto con l'\cref{eq:hall-tension}, che $ q $ sia compatibile con $ 0 $ mentre $ m $ dia una stima di $ R_H B / t $, dopo aver sottratto (da $ m $ e $ q $) il contributo dato dal disallineamento delle piastre.

\subsection{Stima di $ R_H $}

L'andamento teorico dei coefficienti angolari di $ V_H(i_p) $ (debitamente corretti) contro il campo magnetico è lineare.
Per questo effettuiamo una nuova regressione lineare con il modello:
\begin{equation}
m = q' + m' \cdot B
\end{equation}
aspettandoci che $ q' $ sia compatibile con $ 0 $, e che $ m' $ dia una stima di $ R_H / t $.

Sperimentalmente abbiamo misurato la corrente di alimentazione delle bobine del magnete: per ricavare i valori corrispondenti di campo magnetico all'interno del traferro abbiamo usato l'andamento sul ciclo di isteresi medio (\cref{eq:hysteresis-mean}).

Dalla regressione otteniamo \ldots
